This model uses Hugging Face's \textbf{XLNetForSequenceClassification} architecture initialized with the XLNet-base-cased pre-trained checkpoint and tokenizer, with inputs \textbf{padded and truncated to a maximum length of 128 tokens} to ensure consistent dimensions. During training over \textbf{3 epochs} with a \textbf{per-device batch size of 8} (creating an \textbf{effective batch size of 16} via \textbf{2 gradient accumulation steps}), the model utilizes the \textbf{AdamW optimizer}, a \textbf{learning rate of 2e-5}, a \textbf{weight decay of 0.01}, and \textbf{FP16 mixed-precision} to accelerate computation. The task's objective function relies on the standard \textbf{Cross-Entropy Loss} inherent to Hugging Face sequence classification models, and the \textbf{Trainer} is configured to automatically load the \textbf{best-performing weights based on the validation F1-score} at the end of the run.

The model achieved an \textbf{85.21\% accuracy} and an \textbf{85.21\% macro F1-score}, significantly outperforming a \textbf{simple majority-class baseline} (\textbf{F1 = 0.33, Acc = 0.50}) but performing very similar to the \textbf{TF--IDF Logistic Regression baseline} (\textbf{F1 = 0.85, Acc = 0.85}). This similarity in performance is notable: while \textbf{XLNet uses permutation language modeling} to capture \textbf{deep bidirectional semantic context}, TF--IDF relies strictly on \textbf{surface-level word frequency patterns}. The lack of a significant performance gap suggests that the genres in this specific dataset are highly distinguishable through \textbf{surface-level vocabulary cues}, where TF--IDF excels, rather than requiring \textbf{complex contextual semantic understanding}. Alternatively, \textbf{increasing the sequence length beyond 128 tokens} or conducting further \textbf{hyperparameter tuning} might be necessary for the model to further improve its performance.